% !TEX program = pdflatex
% TEHORA — Développeur portail / Tableaux de bord — Jumeau numérique · FR
% Honnête : React, Node, REST, WebRTC/temps réel, dashboards ; junior ≠ 3–7 ans
\documentclass[10pt,a4paper]{article}

\usepackage[french]{babel}
\usepackage[utf8]{inputenc}
\usepackage[T1]{fontenc}
\usepackage{lmodern}
\usepackage[margin=1.25cm]{geometry}
\usepackage{xcolor}
\usepackage{titlesec}
\usepackage{enumitem}
\usepackage{hyperref}
\usepackage{fontawesome5}
\usepackage{microtype}

\setlength{\parindent}{0pt}
\setlength{\parskip}{2pt}
\pagestyle{empty}

\definecolor{primary}{HTML}{1A365D}
\definecolor{accent}{HTML}{2B6CB0}
\definecolor{muted}{HTML}{4A5568}
\definecolor{rulecolor}{HTML}{E2E8F0}

\newcommand{\TargetTitle}{Développeur Front-End / Portails \& Tableaux de bord \textbar{} React}
\newcommand{\TargetKeywords}{React · Node.js · API REST · Tableaux de bord · Visualisation · WebSocket/temps réel · UX · Tailwind · Git · Docker · CI/CD}

\hypersetup{
  colorlinks=true,
  linkcolor=accent,
  urlcolor=accent,
  pdfauthor={Hamouda Chekir},
  pdftitle={CV — Hamouda Chekir — TEHORA Portail / Dashboards}
}

\titleformat{\section}
  {\color{primary}\large\bfseries}
  {}
  {0pt}
  {}
  [\vspace{1pt}{\color{accent}\titlerule[1.6pt]}\vspace{5pt}]
\titlespacing*{\section}{0pt}{10pt}{3pt}

\newcommand{\job}[3]{%
  \textbf{\color{primary}#1}\hfill{\color{muted}\small #2}\\[-1pt]
  {\color{accent}\textit{#3}}\\[2pt]
  \begin{itemize}[leftmargin=0.85em, itemsep=1pt, topsep=0pt, parsep=0pt, partopsep=0pt]
}

\newcommand{\proj}[3]{%
  \textbf{\color{primary}#1}\hfill{\color{muted}\small #2}\\[-1pt]
  {\color{accent}\textit{#3}}\\[1pt]
  \begin{itemize}[leftmargin=0.85em, itemsep=1pt, topsep=0pt, parsep=0pt, partopsep=0pt]
}

\newcommand{\edu}[3]{%
  \textbf{\color{primary}#1}\hfill{\color{muted}\small #2}\\[-1pt]
  {\color{accent}\textit{#3}}\\[3pt]
}

\newcommand{\cert}[4]{%
  \textbf{\color{primary}\small #1}\hfill{\color{muted}\scriptsize #3}\\[-1pt]
  {\color{accent}\small\textit{#2}}\\[0pt]
  {\color{muted}\scriptsize #4}\\[4pt]
}

\newcommand{\contactsep}{\hspace{5pt}{\color{rulecolor}|}\hspace{5pt}}

\begin{document}

{\fontsize{22}{26}\selectfont\bfseries\color{primary}HAMOUDA CHEKIR}\\[2pt]
{\large\color{accent}\TargetTitle}\\[4pt]
{\footnotesize\color{muted}%
  \faIcon{envelope}\hspace{3pt}\href{mailto:hamoudachkir@yahoo.fr}{hamoudachkir@yahoo.fr}%
  \contactsep
  \faIcon{phone}\hspace{3pt}+216\,55\,913\,832%
  \contactsep
  \faIcon{map-marker-alt}\hspace{3pt}Tunis, Tunisie \textbar{} hybride / remote%
}\\[2pt]
{\footnotesize\color{muted}%
  \faIcon{linkedin}\hspace{3pt}\href{https://www.linkedin.com/in/hamouda-chekir-5053572b4}{hamouda-chekir}%
  \contactsep
  \faIcon{github}\hspace{3pt}\href{https://github.com/hamoudachekir}{hamoudachekir}%
}\\[6pt]

\section{PROFIL}
Développeur junior (ESPRIT --- spécialité TWIN : Technologies du Web, de l'Internet et des Réseaux), orienté \textbf{interfaces web, portails et tableaux de bord}.
À l'aise en React et Node.js, consommation d'API REST, parcours temps réel et collaboration avec les besoins métier / UX.
Objectif : contribuer à des portails de jumeau numérique en rendant les données exploitables (vues, alertes, exports) avec performance et clarté d'usage.
Français courant ; anglais professionnel.

\section{EXPÉRIENCE}

\job{Ingénieur Logiciel — Projet de Fin d'Études}{Févr. 2026 -- Août 2026}{Talan \textbar{} Tunis}
  \item Contribué aux \textbf{2 SPA React} de NextHire (parcours candidat / recruteur) : interfaces maintenables, consommation d'API REST et attention portée à l'ergonomie.
  \item Intégré des retours \textbf{temps réel} (sessions WebRTC) avec indicateurs d'état et feedback utilisateur --- logique proche des dashboards live.
  \item Participé à des écrans de reporting / scoring pour les recruteurs (données agrégées, priorisation visuelle) en lien avec les services data / back-end.
  \item Documenté et livré via GitLab CI/CD et Docker, en collaboration multi-équipes (front, back, produit).
\end{itemize}

\job{Stagiaire Développeur Full-Stack}{Juin 2025 -- Août 2025}{SmartConseil — Dafe Management \textbar{} Tunis}
  \item Développé des interfaces back-office (filtres, listes, exports PDF/Excel) orientées opérateurs --- type portail métier / tableau de suivi.
  \item Travaillé sur la clarté UX et la maintenabilité, avec revues de code, tests et CI/CD.
\end{itemize}

\job{Stagiaire Développeur Web}{Juin 2024 -- Août 2024}{Big Deal (Club Privilèges) \textbar{} Tunis}
  \item Livré un module métier (codes cadeaux) avec écrans de gestion et validation avant production.
\end{itemize}

\job{Stagiaire Développeur Web}{Févr. 2023 -- Mai 2023}{HOTIX \textbar{} Tunis}
  \item Mis en production un storefront e-commerce (.NET) sous Scrum, avec focus expérience utilisateur.
\end{itemize}

\section{PROJETS}
{\footnotesize
\setlength{\parskip}{1pt}
\proj{TalentMatch}{2025}{React · FastAPI · visualisation / UX produit}
  \item Front React d'une plateforme d'enchères : vues métier, intégration API et lisibilité des informations clés.
\end{itemize}
\proj{TuniVert}{2025}{Laravel · JS · dashboard admin}
  \item Dashboard admin (analytics, CRUD) et parcours utilisateur --- pratique des tableaux de bord et exports.
\end{itemize}
\proj{Full Stack JS (MERN)}{2023}{React · Node.js · REST}
  \item SPA React + API Node, auth et rôles --- architecture web moderne.
\end{itemize}
}

\section{FORMATION}
\edu{Cycle Ingénieur — Spécialité TWIN (Technologies du Web, de l'Internet et des Réseaux)}{2023 -- 2026}{ESPRIT \textbar{} Tunis}
\edu{Licence Business Computing — E-Business}{2020 -- 2023}{ESEN \textbar{} Manouba}

\section{COMPÉTENCES}
{\small
\textbf{\color{primary}Front / portails :} React, JavaScript, HTML/CSS, Tailwind, composants UI, gestion d'état (bases), ergonomie\\[2pt]
\textbf{\color{primary}Données \& APIs :} API REST, Node/Express, Python FastAPI, PostgreSQL/MongoDB ; intérêt visualisation (Chart.js / ECharts à approfondir)\\[2pt]
\textbf{\color{primary}Temps réel :} WebRTC / flux live, retours d'état utilisateur (base pour WebSocket / streaming)\\[2pt]
\textbf{\color{primary}Livraison :} Git, GitLab CI/CD, Docker, documentation, tests, Agile\\[2pt]
\textbf{\color{primary}Cloud :} AWS Cloud Foundations\\[2pt]
\textbf{\color{primary}Pour ce poste :} \TargetKeywords
}

\vspace{2pt}
\begin{minipage}[t]{0.65\textwidth}
\section{CERTIFICATIONS}
\cert{AWS Academy — Cloud Foundations}{AWS Academy}{Mai 2025}{Socle cloud (AWS)}
\end{minipage}%
\hfill
\begin{minipage}[t]{0.32\textwidth}
\section{LANGUES}
{\small\textbf{Arabe} — Maternel}\\[2pt]
{\small\textbf{Français} — Courant}\\[2pt]
{\small\textbf{Anglais} — Professionnel}
\end{minipage}

\end{document}
